\documentclass[12pt]{article}
\usepackage[utf8]{inputenc}
\usepackage{amsmath,amssymb,amsthm,algorithm,algorithmic}
\usepackage[margin=1in]{geometry}

\newtheorem{theorem}{Theorem}
\newtheorem{lemma}[theorem]{Lemma}
\newtheorem{proposition}[theorem]{Proposition}
\theoremstyle{definition}
\newtheorem{definition}[theorem]{Definition}

\title{Problem 31: Coin Sums}
\author{}
\date{}

\begin{document}
\maketitle

\section{Problem Statement}
In England the currency is made up of pound and pence, with eight coins in general circulation:
\[
1\text{p},\ 2\text{p},\ 5\text{p},\ 10\text{p},\ 20\text{p},\ 50\text{p},\ 100\text{p}\ (\pounds 1),\ 200\text{p}\ (\pounds 2).
\]
How many different ways can $\pounds 2$ (200p) be made using any number of these coins?

\section{Mathematical Development}

\begin{definition}
Let $S = \{c_1, c_2, \ldots, c_k\}$ be a finite set of positive integers (coin denominations). A \emph{representation} of a non-negative integer $n$ over $S$ is a tuple $(j_1, j_2, \ldots, j_k) \in \mathbb{Z}_{\ge 0}^k$ satisfying $\sum_{i=1}^{k} j_i c_i = n$. The \emph{change-making function} $p(n; S)$ counts the number of such representations.
\end{definition}

\begin{theorem}[Generating function]
The change-making function satisfies
\[
p(n; S) = [x^n] \prod_{i=1}^{k} \frac{1}{1 - x^{c_i}},
\]
where $[x^n]$ denotes extraction of the coefficient of $x^n$.
\end{theorem}

\begin{proof}
Each factor admits the geometric series expansion $\frac{1}{1 - x^{c_i}} = \sum_{j=0}^{\infty} x^{j c_i}$, valid for $|x| < 1$. The Cauchy product of $k$ such series yields
\[
\prod_{i=1}^{k} \frac{1}{1 - x^{c_i}} = \sum_{(j_1, \ldots, j_k) \in \mathbb{Z}_{\ge 0}^k} x^{j_1 c_1 + j_2 c_2 + \cdots + j_k c_k}.
\]
The coefficient of $x^n$ equals the cardinality of $\{(j_1, \ldots, j_k) \in \mathbb{Z}_{\ge 0}^k : \sum_{i=1}^{k} j_i c_i = n\}$, which is $p(n; S)$ by definition.
\end{proof}

\begin{lemma}[Recurrence]
Define $f(i, m)$ as the number of representations of $m$ using only denominations $\{c_1, \ldots, c_i\}$. Then:
\[
f(i, m) = \begin{cases}
1 & \text{if } m = 0, \\
0 & \text{if } m > 0 \text{ and } i = 0, \\
f(i-1, m) & \text{if } 0 < m < c_i, \\
f(i-1, m) + f(i, m - c_i) & \text{if } m \ge c_i.
\end{cases}
\]
\end{lemma}

\begin{proof}
The case $m = 0$ admits exactly the empty sum. The case $i = 0$, $m > 0$ has no coins available. For $m \ge c_i$, the representations partition into two disjoint classes:
\begin{enumerate}
\item Those with $j_i = 0$: counted by $f(i-1, m)$.
\item Those with $j_i \ge 1$: subtracting one copy of $c_i$ gives a bijection with representations of $m - c_i$ over $\{c_1, \ldots, c_i\}$, counted by $f(i, m - c_i)$.
\end{enumerate}
These classes are exhaustive and mutually exclusive. For $0 < m < c_i$, we must have $j_i = 0$, collapsing to $f(i-1, m)$.
\end{proof}

\begin{theorem}[Space-optimal computation]
The function $f(k, n)$ can be computed in $O(kn)$ time and $O(n)$ space using a one-dimensional array $\mathrm{dp}[0..n]$.
\end{theorem}

\begin{proof}
Initialize $\mathrm{dp}[0] = 1$ and $\mathrm{dp}[m] = 0$ for $1 \le m \le n$, encoding $f(0, m)$.

\textbf{Claim.} After processing coin $c_i$ via the update $\mathrm{dp}[j] \mathrel{+}= \mathrm{dp}[j - c_i]$ for $j = c_i, c_i + 1, \ldots, n$ in increasing order, we have $\mathrm{dp}[m] = f(i, m)$ for all $0 \le m \le n$.

We prove the claim by induction on $i$. The base case $i = 0$ holds by initialization. Suppose $\mathrm{dp}[m] = f(i-1, m)$ at the start of iteration $i$. For $j \ge c_i$: before $j$ is processed, $\mathrm{dp}[j] = f(i-1, j)$; since $j - c_i < j$ and we process in increasing order, $\mathrm{dp}[j - c_i]$ has already been updated to $f(i, j - c_i)$; after the update, $\mathrm{dp}[j] = f(i-1, j) + f(i, j - c_i) = f(i, j)$ by the lemma. For $j < c_i$, no update occurs, and $\mathrm{dp}[j] = f(i-1, j) = f(i, j)$. After all $k$ iterations, $\mathrm{dp}[n] = f(k, n) = p(n; S)$.
\end{proof}

\section{Algorithm}

\begin{algorithm}[H]
\caption{Coin Change Count}
\begin{algorithmic}[1]
\REQUIRE Target amount $n$, coin set $\{c_1, \ldots, c_k\}$
\ENSURE $p(n; S)$
\STATE $\mathrm{dp}[0..n] \gets 0$; \quad $\mathrm{dp}[0] \gets 1$
\FOR{each $c$ in coins}
    \FOR{$j \gets c$ \TO $n$}
        \STATE $\mathrm{dp}[j] \gets \mathrm{dp}[j] + \mathrm{dp}[j - c]$
    \ENDFOR
\ENDFOR
\RETURN $\mathrm{dp}[n]$
\end{algorithmic}
\end{algorithm}

\section{Complexity Analysis}

\begin{proposition}
For $S = \{1, 2, 5, 10, 20, 50, 100, 200\}$ and $n = 200$: the algorithm performs $\Theta(kn) = \Theta(1600)$ arithmetic operations using $\Theta(n) = \Theta(201)$ integers of storage.
\end{proposition}

\begin{proof}
Follows from Theorem~2 with $k = 8$ and $n = 200$. Each operation is $O(1)$ since all intermediate values fit in a machine word ($f(8,200) = 73682 < 2^{31}$).
\end{proof}

\section{Answer}

\[
\boxed{73682}
\]

\end{document}
